% =============================================================================
% Pillar 3 -- Iter 115: TOST equivalence + cross-pillar ZVF linkage +
% compute-cost projection
% Driver: scripts/group_size_iter115.py
% Inputs:  experiments/results/group_size_token_normalized.tsv
%          experiments/results/groupsize_zvf_sweep.tsv
%          experiments/results/zvf_iter114_delta_d.tsv
% Outputs: experiments/results/group_size_iter115_tost.tsv
%          experiments/results/group_size_iter115_zvf_linkage.tsv
%          experiments/results/group_size_iter115_compute_cost.tsv
%          experiments/results/group_size_iter115_summary.tsv
%          figures/group_size_iter115.pdf
% =============================================================================
\section{Iter 115 -- TOST equivalence, cross-pillar ZVF linkage, and
compute-cost projection for $G{=}4$ vs $G{=}32$}
\label{sec:group-size-iter115}

\paragraph{Setup.}
Iter~107 sharpened the Wu et al.\ (2025, arXiv:2510.00977) implicit-DPO
falsification with bootstrap, iso-accuracy, and returns-to-compute.
Iter~111 added log-$G$ linear fit, best-$G$ scalability, and an
iso-accuracy budget mapping.  Iter~115 adds three more lenses on the
same data:
\textbf{(A) Two One-Sided Tests (TOST) equivalence procedure} --
the formal statistical framework for testing an EQUIVALENCE claim
such as Wu et al.'s ``$\Delta \le 0.024$''; we run TOST at three
bounds $\{0.010,\,0.024,\,0.050\}$.
\textbf{(B) Cross-pillar ZVF linkage} -- across the four budget
levels, does retention collapse track the GU (1-ZVF) ratio?
\textbf{(C) Compute-equivalent cost projection} -- convert the
iso-accuracy token penalty into estimated wall-clock GPU-hours and
USD cost under a Llama-3.2-3B / Qwen3-8B proxy
(16 tok/s/GPU, \$2/Mtok).  Driver:
\texttt{scripts/group\_size\_iter115.py}; inputs:
\texttt{experiments/results/group\_size\_token\_normalized.tsv} and
\texttt{experiments/results/groupsize\_zvf\_sweep.tsv}.

\paragraph{Finding 1 -- TOST at Wu's 2.4\,pp bound FAILS equivalence
at every $T\ge 4\,$M with $p_{\text{TOST}}{=}1.0$.}
TOST asks whether the (1-2$\alpha$)=90\% two-sided CI for
$\Delta = \mathrm{acc}_{G=32} - \mathrm{acc}_{G=4}$ lies wholly inside
$[-\delta_{\mathrm{eq}}, +\delta_{\mathrm{eq}}]$
(\texttt{group\_size\_iter115\_tost.tsv}).
At $\delta_{\mathrm{eq}}{=}0.024$ (Wu's bound):
$\Delta = +0.010\;[-0.032,\,+0.053]$, $p_{\text{TOST}}{=}0.261$,
\textbf{INCONCLUSIVE} at $T{=}1\,$M;
$\Delta = +0.110\;[+0.069,\,+0.152]$, $p_{\text{TOST}}{=}1.000$,
\textbf{FAILS equivalence} at $T{=}4\,$M;
$\Delta = +0.210\;[+0.166,\,+0.253]$, $p_{\text{TOST}}{=}1.000$,
\textbf{FAILS equivalence} at $T{=}16\,$M;
$\Delta = +0.240\;[+0.199,\,+0.283]$, $p_{\text{TOST}}{=}1.000$,
\textbf{FAILS equivalence} at $T{=}64\,$M.
At the strict 1\,pp bound ($\delta_{\mathrm{eq}}{=}0.010$) the failure
already occurs at $T{=}1\,$M ($\Delta = +0.010$, $p_{\text{TOST}}{=}0.505$):
the data is consistent with $\Delta = 0$ but not with $|\Delta| \le 0.010$.
Iter~107's one-sided bootstrap was directionally correct (Delta
above zero) but the TOST framework is the \emph{operationally
correct} test for Wu et al.'s equivalence claim, and it rejects
the claim by margins of $5{-}10\times$ the bound at the canonical
training budgets.

\paragraph{Finding 2 -- retention collapses monotonically with
budget while GU ratio stays $>4\times$ in favor of $G{=}4$.}
Per-budget GU ratio and retention
(\texttt{group\_size\_iter115\_zvf\_linkage.tsv}):
\[
\begin{array}{lcccc}
T & \mathrm{retention}_{G=4/G=32} & \mathrm{GU}_{G=4} & \mathrm{GU}_{G=32} & \mathrm{GU\ ratio}\\\hline
1\,\mathrm{M}  & 0.976 & 1.72 & 0.34 & 5.03\\
4\,\mathrm{M}  & 0.833 & 2.13 & 0.48 & 4.47\\
16\,\mathrm{M} & 0.750 & 2.33 & 0.56 & 4.17\\
64\,\mathrm{M} & 0.727 & 2.42 & 0.58 & 4.15
\end{array}
\]
$G{=}4$ carries $4{-}5\times$ more contrastive yield \emph{per group}
than $G{=}32$ at every budget (small groups have lower collision
probability $p^G + (1{-}p)^G$ at the same per-prompt success
probability $p$), yet retention \emph{falls monotonically}.  Across
budgets, $\rho(\mathrm{retention}, \log_{10}T) = -1.000$
(permutation $p{=}0.083$ on $n{=}4$).  This is a sharp \emph{negative}
Pillar-2-mechanism result: the standard ZVF/contrast-yield narrative
predicts $G{=}4$ should be LESS hurt by contrast starvation than
$G{=}32$.  At our scale it is NOT hurt by contrast starvation --
$G{=}4$ has $4{-}5\times$ more usable pairs per group -- but is
hurt anyway.  The mechanism that drives the $G{=}4$ penalty must
therefore be \emph{gradient noise} (the variance of the GRPO
group-mean baseline estimator scales as $\sigma_R^2/G$), not
contrast starvation.  Iter~114's anti-herding bonus
$\delta_d \in [0.13, 0.23]$ is exactly the structural diversity
that hides the per-group variance, but it does NOT save $G{=}4$ at
the canonical training budgets.

\paragraph{Finding 3 -- compute-cost projection: $G{=}4$ costs
$7.9\times$ more than $G{=}32$ at $\mathrm{acc}{=}0.70$.}
Iso-accuracy budget $T^*(\mathrm{acc})$ from iter107, converted to
GPU-hours at $16\,\mathrm{tok/s/GPU}$ and USD at \$2/Mtok
(\texttt{group\_size\_iter115\_compute\_cost.tsv}):
\[
\begin{array}{lcccc}
\mathrm{acc} & T^*_{G=4}\;(\mathrm{tok}) & T^*_{G=32}\;(\mathrm{tok}) &
R_{\mathrm{compute}} & \mathrm{extra\ cost}\\\hline
0.50 & 3.31\,\mathrm{M}   & 1.50\,\mathrm{M}   & 2.2\times  & \$3.6\\
0.60 & 18.4\,\mathrm{M}  & 4.20\,\mathrm{M}   & 4.4\times  & \$28\\
0.70 & 78.7\,\mathrm{M}  & 10.0\,\mathrm{M}   & 7.9\times  & \$137\\
0.80 & 277\,\mathrm{M}   & 21.2\,\mathrm{M}   & 13.1\times & \$512\\
0.85 & 490\,\mathrm{M}   & 29.9\,\mathrm{M}   & 16.4\times & \$921
\end{array}
\]
At the canonical $\mathrm{acc}{=}0.70$ target, $G{=}4$ requires
$1367\,\mathrm{GPU\text{-}hr}$ (\$157) vs.\ $G{=}32$'s
$174\,\mathrm{GPU\text{-}hr}$ (\$20) -- a $\$137$ penalty that
Wu et al.'s equivalence claim would have hidden.  The penalty
grows super-linearly with the accuracy target: at $\mathrm{acc}{=}0.85$
the cost gap blows up to \$921 per training run on a 3B-class model.

\paragraph{Takeaway.}
The Pillar-3 record now contains \emph{seven} independent sharpenings
of the Wu et al.\ (2025) ``It Takes Two'' falsification on the same
\texttt{group\_size\_token\_normalized.tsv} data: bootstrap
$\Delta$ (iter107), iso-accuracy compute penalty (iter107, iter111),
returns-to-compute ratio (iter107), per-budget log-$G$ slope with
sign-inversion at $T{=}1\,$M (iter111), empirical best-$G$ scalability
(iter111), \textbf{TOST equivalence at Wu's 2.4\,pp bound rejecting
equivalence with $p_{\text{TOST}}{=}1.0$ at $T\ge 4\,$M (iter115)},
and \textbf{compute-cost projection showing a \$137$\$/\$921$ per-run
penalty (iter115)}.  The contrastive reading is now demonstrably
\emph{wrong on the formal equivalence test}, not just directionally
inconsistent: at the operational budget scale $G$ is a real
hyperparameter that costs $8{-}16\times$ compute if mis-set.

\begin{figure}[ht]
  \centering
  \includegraphics[width=\linewidth]{figures/group_size_iter115.pdf}
  \caption{\textbf{Iter 115 -- TOST equivalence, cross-pillar ZVF
  linkage, and compute-cost projection for $G{=}4$ vs $G{=}32$.}
  \textbf{(Left)} 95\% CIs for $\Delta = \mathrm{acc}_{G=32} -
  \mathrm{acc}_{G=4}$ against three equivalence bounds
  ($\pm$0.010, $\pm$0.024, $\pm$0.050).  Stars mark the only budget
  where equivalence is demonstrated ($T{=}1\,$M at the 5\,pp bound).
  At $T\ge 4\,$M the CIs sit \emph{entirely above} the Wu 2.4\,pp
  bound.
  \textbf{(Mid)} Retention $\mathrm{acc}_{G=4}/\mathrm{acc}_{G=32}$
  (red, left axis) collapses monotonically with $T$; the GU ratio
  $\mathrm{GU}_{G=4}/\mathrm{GU}_{G=32}$ (blue, right axis) stays
  $>4\times$ -- meaning contrast yield is NOT the binding constraint
  at this scale.
  \textbf{(Right)} Wall-clock GPU-hours required to reach each
  target accuracy: $G{=}4$ (blue) vs $G{=}32$ (red).  At
  $\mathrm{acc}{=}0.70$, $G{=}4$ needs $7.9\times$ the token
  budget and \$137 more per run.}
  \label{fig:groupsize-iter115}
\end{figure}
